\documentclass[11pt]{article}
\usepackage[margin=1in]{geometry}
\usepackage{enumitem}
% =============================================================================
% Preambles/header.tex — shared LaTeX/Beamer preamble for this project
%
% Use from a Beamer lecture by prepending:
%     \input{header}
% and compiling with TEXINPUTS=../Preambles:$TEXINPUTS (the /compile-latex
% skill does this automatically).
%
% PALETTE CONTRACT. Color names defined here MUST match the names in
% Quarto/theme-template.scss so Beamer and Quarto renderings use the same
% palette. The scripts/check-palette-sync.sh script enforces this.
%
% To customize: edit the HEX values below AND the matching $variable in
% Quarto/theme-template.scss, then run scripts/check-palette-sync.sh.
% =============================================================================

% -----------------------------------------------------------------------------
% Engine detection + encoding
%
% Our /compile-latex skill uses XeLaTeX, where inputenc is unnecessary and
% can produce encoding warnings. Only load inputenc under pdfTeX.
% -----------------------------------------------------------------------------
\usepackage{iftex}
\ifPDFTeX
  \usepackage[utf8]{inputenc}
\fi

\usepackage{amsmath, amssymb, amsthm}
\usepackage{booktabs}
\usepackage{graphicx}

% -----------------------------------------------------------------------------
% PALETTE — mirrors Quarto/theme-template.scss variable names.
% Load xcolor and define colors BEFORE anything that references them
% (notably hyperref, hypersetup, and the Beamer color assignments).
% -----------------------------------------------------------------------------
\usepackage{xcolor}

\definecolor{primary-blue}{HTML}{012169}      % headings, accents
\definecolor{primary-gold}{HTML}{B9975B}      % emphasis, borders
\definecolor{highlight-yellow}{HTML}{F2A900}  % markers, alerts
\definecolor{light-bg}{HTML}{E8EDF5}          % title / section backgrounds
\definecolor{jet}{HTML}{1A1A1A}               % body text

% Semantic colors (used in diagrams and annotations)
\definecolor{positive}{HTML}{15803D}          % good / observed / identified
\definecolor{negative}{HTML}{B91C1C}          % bad / problematic / bias
\definecolor{neutral}{HTML}{525252}           % reference / context

% Highlight accents (match .hi-* classes in the SCSS)
\definecolor{hi-slate}{HTML}{314F4F}
\definecolor{hi-green}{HTML}{2E8B57}
\definecolor{hi-red}{HTML}{C41E3A}

% Semantic aliases so skill templates and snippets can write
% \textcolor{accent}{...} without hard-coding "primary-blue".
\colorlet{accent}{primary-blue}
\colorlet{accent2}{primary-gold}

% -----------------------------------------------------------------------------
% Hyperref — loaded AFTER xcolor + palette so link colors resolve.
% -----------------------------------------------------------------------------
\usepackage{hyperref}
\hypersetup{colorlinks=true, linkcolor=primary-blue, urlcolor=primary-blue,
            citecolor=primary-blue, pdfborder={0 0 0}}

% -----------------------------------------------------------------------------
% Course code — set once per deck (after \input{header}, before \begin{document})
% via \coursecode{CS401}. Shown in the footer on every slide so a deck's course
% is identifiable without opening the title page. Empty by default.
% -----------------------------------------------------------------------------
\makeatletter
\newcommand{\coursecode}[1]{\gdef\@coursecodetext{#1}}
\gdef\@coursecodetext{}
\makeatother

% -----------------------------------------------------------------------------
% Beamer theme assignments (only apply under Beamer).
%
% We detect Beamer via \@ifclassloaded{beamer}, which is the documented,
% non-internal way. This must be wrapped in \makeatletter/\makeatother
% because \@ifclassloaded contains an @-catcode character.
% -----------------------------------------------------------------------------
\makeatletter
\@ifclassloaded{beamer}{%
  \usefonttheme{professionalfonts}
  \setbeamercolor{structure}{fg=primary-blue}
  \setbeamercolor{frametitle}{fg=primary-blue}
  \setbeamercolor{title}{fg=primary-blue}
  \setbeamercolor{subtitle}{fg=primary-gold}
  \setbeamercolor{author}{fg=primary-blue}
  \setbeamercolor{institute}{fg=neutral}
  \setbeamercolor{date}{fg=primary-gold}
  \setbeamercolor{itemize item}{fg=primary-blue}
  \setbeamercolor{itemize subitem}{fg=primary-gold}
  \setbeamercolor{alerted text}{fg=primary-gold}
  \setbeamercolor{block title}{fg=white, bg=primary-blue}
  \setbeamercolor{block body}{bg=light-bg}
  % Semantic block variants: block=definition/concept (blue), exampleblock=
  % worked example (green/positive), alertblock=key takeaway or caution (gold).
  % Keep to <=2 colored boxes per slide across ALL three variants (INV-7).
  \setbeamercolor{block title example}{fg=white, bg=positive}
  \setbeamercolor{block body example}{bg=light-bg}
  \setbeamercolor{block title alerted}{fg=white, bg=primary-gold}
  \setbeamercolor{block body alerted}{bg=light-bg}

  % Minimal footer: course code (left) + page X / N (right), no navigation clutter
  \setbeamertemplate{navigation symbols}{}
  \setbeamertemplate{footline}{%
    \color{neutral}\footnotesize\hspace{0.7em}\@coursecodetext\hfill
    \insertframenumber/\inserttotalframenumber\hspace{0.7em}\vspace{0.3em}}
}{}
\makeatother

% -----------------------------------------------------------------------------
% TikZ setup — libraries the snippet gallery uses
% -----------------------------------------------------------------------------
\usepackage{tikz}
\usetikzlibrary{arrows.meta, positioning, calc, decorations.pathreplacing,
                fit, shapes.geometric, backgrounds}

% Shared TikZ styles — snippets and hand-written diagrams can pick these up
% via `\begin{tikzpicture}[every node/.style={...}]` or by naming specific
% styles (e.g., `\node[dag-node] (X) at (0,0) {X};`).
\tikzset{
  % Node styles for causal DAGs and similar
  dag-node/.style = {
    draw=primary-blue, fill=light-bg, rounded corners=2pt,
    minimum width=1.2cm, minimum height=0.9cm, align=center, font=\small
  },
  decision-node/.style = {
    draw=primary-gold, fill=white, diamond, aspect=1.8,
    minimum width=1.8cm, minimum height=1.2cm, align=center, font=\small,
    inner sep=1pt
  },
  % Edge styles — observed vs counterfactual
  observed-edge/.style = {-{Latex[length=2.5mm]}, thick, draw=jet},
  counterfactual-edge/.style = {-{Latex[length=2.5mm]}, thick, draw=neutral, dashed},
  confound-edge/.style = {-{Latex[length=2.5mm]}, thick, draw=negative, dashed},
  % Data-point styles for plots
  observed-dot/.style = {fill=primary-blue, draw=primary-blue, circle, inner sep=1.2pt},
  counterfactual-dot/.style = {fill=white, draw=neutral, circle, inner sep=1.2pt}
}

% -----------------------------------------------------------------------------
% Convenience macros
% -----------------------------------------------------------------------------
\newcommand{\muted}[1]{\textcolor{neutral}{#1}}
\newcommand{\key}[1]{\textcolor{primary-gold}{\textbf{#1}}}
\newcommand{\good}[1]{\textcolor{positive}{#1}}
\newcommand{\bad}[1]{\textcolor{negative}{#1}}

% Standout frame macro: full-bleed dark background for section transitions.
% Beamer-only (uses \begin{frame}/\setbeamercolor/\usebeamerfont) -- do not
% call from an article-class file (e.g. Notes/<CODE>/*-notes.tex). \muted,
% \key, \good, \bad, and \coursecode above are plain xcolor/gdef and are
% safe to use from either class; verified when Notes/article-class support
% was added.
% Expand: \transitionslide{Your transition title}
\newcommand{\transitionslide}[1]{%
  \begingroup
  \setbeamercolor{background canvas}{bg=primary-blue}
  \begin{frame}[plain]
    \centering
    \vfill
    {\usebeamerfont{title}\color{white}\Large #1\par}
    \vfill
  \end{frame}
  \endgroup
}

% Section divider frame (Beamer-only), an alternative to \transitionslide with a
% darker two-line style: near-black (jet) background, a small white label line
% above a larger gold title, and a thin gold rule along the bottom edge.
% Expand: \sectiondivider{Section 2}{Pointers}
\newcommand{\sectiondivider}[2]{%
  \begingroup
  \setbeamercolor{background canvas}{bg=jet}
  \begin{frame}[plain]
    \vspace*{3.0cm}
    {\color{white}\ttfamily\large #1\par}
    \vspace{0.5cm}
    {\color{primary-gold}\LARGE #2\par}
    \begin{tikzpicture}[remember picture, overlay]
      \draw[primary-gold, line width=2.5pt]
        ($(current page.south west)+(0,0.1)$) -- ($(current page.south east)+(0,0.1)$);
    \end{tikzpicture}
  \end{frame}
  \endgroup
}

% -----------------------------------------------------------------------------
% Channel watermark — "Concepts That Click" (YouTube).
%
% Article-class files (CompetitiveExam Q/A sets, Notes, Assignments) get a
% light diagonal draft watermark on every page plus a centered footer naming
% the channel. Beamer decks get a subtle TikZ background watermark instead
% (the footline already carries the course code). Centralized here so every
% MCQ PDF is branded without per-file edits.
% -----------------------------------------------------------------------------
\makeatletter
\@ifclassloaded{article}{%
  \usepackage{draftwatermark}
  \SetWatermarkText{Concepts That Click}
  \SetWatermarkScale{1}
  \SetWatermarkFontSize{2cm}
  \SetWatermarkLightness{0.86}
  \SetWatermarkAngle{45}
  \usepackage{fancyhdr}
  \pagestyle{fancy}
  \fancyhf{}
  \fancyfoot[C]{\footnotesize\textcolor{neutral}{Concepts That Click~|~YouTube: \href{https://www.youtube.com/@concepts.thatclick}{@concepts.thatclick}}}
  \renewcommand{\headrulewidth}{0pt}
  \renewcommand{\footrulewidth}{0pt}
  \fancypagestyle{plain}{%
    \fancyhf{}%
    \fancyfoot[C]{\footnotesize\textcolor{neutral}{Concepts That Click~|~YouTube: \href{https://www.youtube.com/@concepts.thatclick}{@concepts.thatclick}}}%
    \renewcommand{\headrulewidth}{0pt}%
    \renewcommand{\footrulewidth}{0pt}%
  }
}{}
\@ifclassloaded{beamer}{%
  \setbeamertemplate{background}{%
    \begin{tikzpicture}[remember picture, overlay]
      \node[rotate=45, opacity=0.055, align=center]
        at (current page.center)
        {\Huge\textcolor{neutral}{Concepts That Click}};
    \end{tikzpicture}%
  }
}{}
\makeatother 
\coursecode{JKSSB}
\renewcommand{\thesection}{34.\arabic{section}}
\newtheorem{example}{Example}[section]
\newenvironment{definitionbox}[1]{%
  \par\noindent\textbf{Definition (#1).}\ \itshape
}{\par\vspace{0.3em}}
\title{Internet, WWW and Browser --- Detailed Lecture Notes}
\author{JKSSB Computer Awareness}
\date{\today}

\begin{document}
\maketitle

\section{The Internet: the network itself}
The Internet is a global network of connected computers and networks. It is
best described as a \textbf{network of networks}: smaller networks joined
together so that any connected device can reach any other. It carries many
services --- the World Wide Web, e-mail, file transfer, and more. The
Internet is the physical infrastructure underneath all of these: the cables,
routers, servers, and the devices connected to them.

\begin{definitionbox}{Internet}
The global network of interconnected computers and networks over which
services such as the Web and e-mail run.
\end{definitionbox}

\begin{center}
\begin{tabular}{p{0.24\textwidth}p{0.66\textwidth}}
\toprule
\textbf{Term} & \textbf{Meaning in this lesson} \\
\midrule
Internet & The global network of networks; the infrastructure. \\
WWW (World Wide Web) & A service on the Internet: linked web pages and sites. \\
Website & A collection of related webpages under one address. \\
Webpage & A single document on the Web, shown in a browser. \\
Homepage & The main or starting page of a website. \\
Browser & Software that opens and displays websites and webpages. \\
Search engine & A website that finds other pages matching typed keywords. \\
Client & The device or program that requests a service or resource. \\
Server & The computer or program that supplies services to clients. \\
ISP (Internet Service Provider) & The company or organisation that provides the Internet connection. \\
Hyperlink & Clickable text or image that jumps to another page or section. \\
URL (Uniform Resource Locator) & The address of a webpage or resource on the Web. \\
\bottomrule
\end{tabular}
\end{center}

The rest of this lesson uses these terms in one consistent sense. The
\textbf{Internet} is the network. The \textbf{WWW} is the linked-pages
service running on that network. A \textbf{webpage} is one document, a
\textbf{website} is the collection, and the \textbf{homepage} is the
collection's front door. A \textbf{browser} opens pages; a \textbf{search
engine} finds pages. A \textbf{client} requests and a \textbf{server}
supplies. An \textbf{ISP} provides the connection itself. A
\textbf{hyperlink} points to another page, and a \textbf{URL} is a page's
address.

\section{Internet vs WWW}
The single most tested confusion in this lesson is Internet versus WWW. The
Internet is the whole global network; the WWW --- World Wide Web --- is one
service that runs on the Internet, namely the system of linked web pages and
sites. E-mail is another service on the same Internet, and it is not part of
the Web. Access follows the same layering: a device reaches the Internet
through an ISP connection, and it reaches the Web through a browser running
over that Internet connection.

\begin{definitionbox}{World Wide Web (WWW)}
The system of linked web pages and websites accessed over the Internet
through a browser; one service among several that the Internet carries.
\end{definitionbox}

\begin{example}
An item states, ``The Internet and the World Wide Web are the same thing.''
This is wrong. The Internet is the global network of networks, while the WWW
is one service running on it. The wrong option has collapsed the
infrastructure and one of its services into a single term.
\end{example}

\section{Website, webpage, and homepage}
A \textbf{webpage} is a single document on the Web, shown inside a browser.
A \textbf{website} is a collection of related webpages grouped under one
address (domain). The \textbf{homepage} is the main or starting page of a
website --- the page that loads when the site's base address is typed. In
order of size: the homepage is one page, any single document is a webpage,
and the whole collection is the website.

\begin{definitionbox}{Website, webpage, homepage}
A webpage is one document on the Web; a website is the collection of related
webpages under one address; the homepage is the website's main or starting
page.
\end{definitionbox}

Each webpage of a website has its own URL, and the homepage URL is the
site's base address. Typing that address in the browser's address bar loads
the homepage first, and hyperlinks on it lead to the site's remaining pages.

\section{Browser}
A \textbf{browser} is software used to open and view websites and webpages.
It sends a request for a page and displays the page that comes back. Common
browsers are Google Chrome, Mozilla Firefox, Microsoft Edge, and Apple
Safari. A browser is software installed on the user's device; it is not the
Internet itself, and it is not a search engine.

\begin{definitionbox}{Browser}
Software that requests webpages and displays them; examples are Google
Chrome, Mozilla Firefox, Microsoft Edge, and Apple Safari.
\end{definitionbox}

\section{Search engine}
A \textbf{search engine} is a website (a service) that finds other webpages
matching keywords typed by the user. Examples are Google, Bing, and Yahoo.
It is reached \emph{through} a browser: the browser is the tool open on the
device, and the search-engine site is loaded inside it. Searching therefore
needs both layers --- a browser open on the device, and a search-engine site
loaded in that browser.

\begin{definitionbox}{Search engine}
A website that returns the pages matching a user's keywords; examples are
Google, Bing, and Yahoo.
\end{definitionbox}

The browser-versus-search-engine table is worth memorising as a whole. A
browser is software (an application) whose role is to open and display pages;
a search engine is a website whose role is to find pages for keywords. Chrome,
Firefox, Edge, and Safari belong in the first column; Google, Bing, and
Yahoo belong in the second. The classic trap swaps them: it calls Chrome a
search engine or Google a browser.

\section{Client and server}
A \textbf{client} is the device or program that requests a service or
resource. The office PC's browser asking for a webpage is acting as a
client: it asks, and it does not supply the page itself. A \textbf{server}
is the computer (or program) that supplies services or resources to clients.
A web server stores websites and sends pages to the browsers that ask for
them. One server can serve many clients at the same time, and servers stay
switched on and connected so that clients can reach them at any time.

\begin{definitionbox}{Client and server}
The client requests a service or resource; the server supplies it. On the
Web, the browser is the client and the web server holds the pages.
\end{definitionbox}

The two always appear as a pair, in three steps. First, the client (the
browser) sends a request for a page. Second, the server receives the request
and sends the page back. Third, the browser displays the received page. This
request-and-response pair is the client-server model of the Web: typing a
site address and pressing Enter sends a request, the site's server returns
the homepage, and the browser shows it.

\section{ISP: the connection provider}
\textbf{ISP} stands for Internet Service Provider. It is the company or
organisation that provides the Internet connection to homes and offices.
Without an ISP connection, the device cannot reach the Internet at all, so
the ISP link comes before everything else: no ISP means no browsing, no
search, and no e-mail. This is the tested fact of WO 2026 Q74 (PYQ-34): ISP
expands to Internet Service Provider, and no other expansion is accepted.

\begin{definitionbox}{ISP (Internet Service Provider)}
The company or organisation that provides the Internet connection to homes
and offices.
\end{definitionbox}

\section{Hyperlink and URL}
A \textbf{hyperlink} (link) is clickable text or an image that jumps to
another page or section. Clicking it loads the linked webpage, which may sit
on the same site or on a different one. Links are what make the Web a
``web'': pages joined to pages. Note the division of labour: a hyperlink
points the way, and the browser follows it when clicked.

A \textbf{URL} --- Uniform Resource Locator --- is the address of a webpage
or resource on the Web. Typing the site's URL in the browser's address bar
loads its homepage. The URL names the page; the hyperlink jumps to it.

\begin{definitionbox}{Hyperlink and URL}
A hyperlink is clickable text or an image that leads to another page; a URL
(Uniform Resource Locator) is the address of a webpage or resource.
\end{definitionbox}

\section{Full-form ladder of this lesson}
Direct-recall full-form items are among the most common JKSSB computer
formats, so learn each expansion exactly and give the full form on first
mention. \textbf{ISP} --- Internet Service Provider --- gives the Internet
connection. \textbf{WWW} --- World Wide Web --- is the linked system of
pages and sites. \textbf{URL} --- Uniform Resource Locator --- is the address
of a page or resource. \textbf{HTTP} --- HyperText Transfer Protocol --- is
the set of rules browsers and servers follow when requesting and delivering
pages. \textbf{HTML} --- HyperText Markup Language --- is the language in
which webpages are written.

\section{Exam retrieval and common confusions}
Five distinctions prevent most errors on this topic.
\begin{enumerate}
  \item The Internet is \textbf{not} the Web; the WWW is \textbf{one
    service} running on the Internet.
  \item A browser is \textbf{not} a search engine. Chrome is a
    \textbf{browser}; Google is a \textbf{search engine}.
  \item A website is the \textbf{collection}; a webpage is \textbf{one
    document}; the homepage is the \textbf{first page}.
  \item The client \textbf{requests} and the server \textbf{supplies}. The
    ISP --- Internet Service Provider --- gives the connection (PYQ-34).
  \item A hyperlink \textbf{points}; the browser \textbf{opens}. A URL is the
    page's address.
\end{enumerate}

\begin{example}
An item asks which software is used to view websites. Trace the roles: the
browser opens and displays pages, while the search engine finds pages and
the ISP provides the connection. The answer is the browser (for example
Chrome). If the stem instead asks which site finds pages for keywords, the
answer is the search engine (for example Google), and if it asks who
provides the connection, the answer is the ISP.
\end{example}

\subsection*{Exercises}
\begin{enumerate}[label=\arabic*.]
  \item Distinguish the Internet from the World Wide Web.
  \item Distinguish a website from a webpage and from a homepage.
  \item Distinguish a browser from a search engine, giving two examples of
    each.
  \item Explain the roles of client, server, and ISP when a page is opened.
  \item Define hyperlink and URL, and explain how the two work together.
\end{enumerate}

\subsection*{Solutions}
\begingroup\small
\begin{enumerate}[label=\arabic*.]
  \item The Internet is the global network of networks --- the
    infrastructure. The WWW (World Wide Web) is one service running on that
    infrastructure: the system of linked web pages and sites.
  \item A webpage is one document on the Web; a website is the collection of
    related webpages under one address; the homepage is the website's main or
    starting page.
  \item A browser is software that opens and displays pages (Chrome, Firefox,
    Edge, Safari). A search engine is a website that finds pages for keywords
    (Google, Bing, Yahoo); it is reached through a browser.
  \item The client (browser) requests the page, the server (web server)
    supplies it, and the ISP (Internet Service Provider) provides the
    Internet connection over which the request and response travel.
  \item A hyperlink is clickable text or an image that jumps to another page;
    a URL (Uniform Resource Locator) is a page's address. Clicking the
    hyperlink tells the browser to load the page named by that URL.
\end{enumerate}
\endgroup
\end{document}
